\documentclass[11pt]{book}

% ---------------------------------------------------------------
% Fonts (engine-agnostic via fontspec; compile with LuaLaTeX)
% ---------------------------------------------------------------
\usepackage{fontspec}
\setmainfont{TeX Gyre Pagella}
% unicode-math would give a closer text/math match but its dependency
% lualatex-math is not present in this build environment; mathpazo is the
% standard NFSS-based Palatino-compatible math companion and needs no
% CTAN fetch.
\usepackage{mathpazo}

% ---------------------------------------------------------------
% Core packages
% ---------------------------------------------------------------
\usepackage[a4paper,margin=1in]{geometry}
\usepackage{amsthm}
\usepackage{mathtools}
\usepackage{amssymb}
\usepackage{physics}
\usepackage{graphicx}
\usepackage{hyperref}
\usepackage[capitalize]{cleveref}
\usepackage{tikz}
\usepackage{tikz-cd}
\usepackage{booktabs}
\usepackage{enumitem}
\usepackage{makeidx}
\makeindex
\usepackage{xcolor}
\usepackage{mdframed}

% A visually distinct callout for explicit misreading warnings -- used at
% every point in the book where a result could plausibly be compressed into
% an unconditional claim it does not make. Deliberately styled differently
% from ordinary remarks so it survives skimming.
\newmdenv[
  linecolor=black,
  linewidth=0.75pt,
  backgroundcolor=black!5,
  leftmargin=0pt,
  rightmargin=0pt,
  innertopmargin=6pt,
  innerbottommargin=6pt,
  innerleftmargin=8pt,
  innerrightmargin=8pt,
  skipabove=10pt,
  skipbelow=10pt
]{scopebox}
\newcommand{\scopetitle}[1]{{\bfseries Scope. }#1}

% ---------------------------------------------------------------
% Theorem environments
% ---------------------------------------------------------------
\theoremstyle{plain}
\newtheorem{theorem}{Theorem}[chapter]
\newtheorem{proposition}[theorem]{Proposition}
\newtheorem{lemma}[theorem]{Lemma}
\newtheorem{corollary}[theorem]{Corollary}

\theoremstyle{definition}
\newtheorem{definition}[theorem]{Definition}
\newtheorem{example}[theorem]{Example}
\newtheorem{assumption}[theorem]{Assumption}
\newtheorem{fact}[theorem]{Fact}

\theoremstyle{remark}
\newtheorem{remark}[theorem]{Remark}

% Assumption-ledger environment: tracks what a chapter has and has not established
\newtheorem{ledgerentry}[theorem]{Ledger}

% ---------------------------------------------------------------
% Custom commands
% ---------------------------------------------------------------
\newcommand{\R}{\mathbb{R}}
\newcommand{\C}{\mathbb{C}}
\newcommand{\Z}{\mathbb{Z}}
\newcommand{\Hset}{\mathbb{H}} % quaternions, to avoid clash with Hilbert space H
\newcommand{\Hil}{\mathcal{H}} % Hilbert space
% \dd is provided by the physics package; not redefined here

% ---------------------------------------------------------------
% Title matter
% No exact publication date on the title page (month/year or year only)
% ---------------------------------------------------------------
\title{Quantum Mechanics from Admissible Probabilistic Composition\\
\large A Reconstruction via Unistochastic Matrices}
\author{Flyxion\\ \textit{Independent Researcher}}
\date{2026}

\begin{document}

\frontmatter
\maketitle

\tableofcontents

\chapter*{Preface}
\addcontentsline{toc}{chapter}{Preface}

\begin{scopebox}
\scopetitle{}This book does \emph{not} conclude that complex numbers are
metaphysically necessary, uniquely forced, or the only mathematics
capable of describing quantum phenomena. Every result ruling out $\R$ does
so \emph{conditional on stated assumptions} — a specific admissibility
constraint (Chapters~\ref{ch:context-graphs}--\ref{ch:real-mub-obstruction}),
or an explicitly adopted principle such as local tomography
(Chapter~\ref{ch:local-tomography}, itself flagged throughout as
\emph{assumed}, not derived). Chapter~\ref{ch:final-audit} collects every
such assumption in one place, exactly four of them, and
Chapter~\ref{ch:final-audit}'s \S\ref{sec:recent-literature} and
Chapter~\ref{ch:amplitwist} discuss published 2025--2026 work showing that
changing the assumption about how systems compose reproduces every quantum
prediction using only real numbers. A reader who takes away "this book
proves the universe requires complex numbers" has taken away a claim this
book explicitly argues against.
\end{scopebox}

This book asks how much of quantum mechanics can be recovered from
constraints on admissible probabilistic transformations, rather than
introduced by postulate. The starting point is elementary: a transition
matrix built from squared amplitudes, $B_{ij}=|U_{ij}|^2$, is the
\emph{unistochastic} object at the center of Życzkowski, Kuś, Słomczyński,
and Sommers's~\cite{zyczkowski2003} work and of Barandes's
stochastic--quantum correspondence~\cite{barandes}. The question this
book pursues is what happens if that object is treated not as a curiosity
about quantum probability, but as the thing to be explained: under what
admissibility constraints on probabilistic data does a coherent unitary lift
become necessary, and how much of the standard formalism follows once it
does?

Two methodological commitments run through every chapter.

\paragraph{Reconstruction versus recovery.} A \emph{reconstruction} claim is
proved from constraints already established earlier in this book. A
\emph{recovery} claim shows that a reconstructed structure reproduces
standard quantum mechanics, checked directly against
Sakurai--Napolitano's~\cite{sakurainapolitano} \textit{Modern Quantum
Mechanics} and Gross's~\cite{gross} \textit{Advanced Quantum Mechanics}
lecture notes. (The older Sakurai \textit{Advanced Quantum
Mechanics}~\cite{sakurai} was part of the originally intended reference
corpus but survives only as a scanned, non-text-searchable copy; it is
listed for completeness but was not directly quotable, and no claim in
this book depends on it.) The two are never allowed
to blur into each other: a result that merely matches the textbook formalism
is not thereby derived, and a result derived from this book's admissibility
logic is not thereby the same claim as its textbook counterpart until the
match is checked explicitly.

\paragraph{The assumption ledger.} Each chapter closes with a boxed ledger
stating plainly what was established and what was not — including, where
relevant, exactly which new physical or operational assumption a result
costs. This is not a stylistic tic. Fifty-some pages into a successful
reconstruction it becomes easy to feel that quantum mechanics has been
derived outright; Chapter~\ref{ch:assumption-audit} and
Chapter~\ref{ch:final-audit} exist specifically to check that feeling against
the record, and the per-chapter ledgers are what make that final audit
possible to write honestly rather than from memory.

\bigskip

The book has three parts. Part I asks why bare probability data is
insufficient on its own — first sequentially (composing transformations loses
information under projection) and then contextually (three individually
admissible measurement contexts can fail to admit any single real
realization at once, the sharpest form of which is a minimal graph-theoretic
witness, the triangle). Part II asks what composing separate systems, rather
than composing or contextualizing a single one, forces and distinguishes
among candidate scalar fields — real, complex, quaternionic — culminating in
a single matrix whose entanglement status depends on which field is doing the
composing. Part III returns to the standard formalism and asks how much of
it can be recovered from what Parts I and II established, closing with an
explicit accounting of every assumption the recovery required.

Four such assumptions appear in the whole book, no more and no fewer, and
they are collected in one place in Chapter~\ref{ch:final-audit}. The reader
who wants to know exactly what standard quantum mechanics cost, on this
route, need only turn there.

One methodological habit recurs often enough to name here rather than leave
implicit: whenever an object in this book appears indispensable — a
particular matrix representation, a particular composition rule, a
particular coordinatization — the question asked is always whether it is a
genuine admissibility constraint or merely one convenient way of writing a
deeper invariant structure. Complex numbers themselves are not exempt from
this question; nor is the specific Kronecker-product convention used to
combine systems throughout Parts I and II
(\S\ref{sec:recent-literature} returns to this point directly, in light of
recent literature showing that composition, not the scalar field alone,
carries much of the weight usually attributed to $i$ itself). A
representation earns its place by what it lets the reader compute or check —
not by familiarity, and not by having no visible alternative.

\mainmatter

\part{Single-System Foundations}
\input{chapters/ch01-probability-simplex}
\input{chapters/ch02-bistochastic-birkhoff}
\input{chapters/ch03-orthostochastic-unistochastic}
\input{chapters/ch04-composition-obstruction}
\input{chapters/ch05-context-graphs-cycles}
\input{chapters/ch06-real-mub-obstruction}

\part{Composite Systems and the Scalar Field}
\input{chapters/ch07-minimal-witnesses}
\input{chapters/ch08-local-tomography}
\input{chapters/ch09-tensor-product}
\input{chapters/ch10-positivity}
\input{chapters/ch11-partial-trace}
\input{chapters/ch12-entanglement}
% further chapters TBD: multipartite entanglement, entanglement measures

\part{Recovered Structure}
\input{chapters/ch13-assumption-audit}
\input{chapters/ch14-born-rule}
\input{chapters/ch15-observables}
\input{chapters/ch16-dynamics}
\input{chapters/ch17-povm-dilation}
\input{chapters/ch18-final-audit}
\input{chapters/ch19-amplitwist}
% Part III complete. Part IV decision pending: Berry phase/holonomy,
% spin/symmetry, multipartite entanglement, or higher-d/quaternionic
% minimal witnesses.

\part{Deepening the Reconstruction}
\input{chapters/ch20-higher-d-witnesses}
\input{chapters/ch21-quaternionic}
\input{chapters/ch22-forbidden-subgraph}
\input{chapters/ch23-density-threshold}
\input{chapters/ch24-quaternionic-mub}
% Part IV nearing a natural stopping point given book length. Remaining
% open: explicit quaternionic MUB construction exceeding d+1; translation
% to a graph-realizability witness.

\appendix
\input{chapters/appA-notation}
\input{chapters/appB-quaternions}

\begin{thebibliography}{99}

\bibitem{sakurainapolitano}
J.~J.~Sakurai and J.~Napolitano, \textit{Modern Quantum Mechanics}, 3rd ed. Cambridge University Press.

\bibitem{sakurai}
J.~J.~Sakurai, \textit{Advanced Quantum Mechanics}. Addison-Wesley.

\bibitem{gross}
D.~Gross, \textit{Advanced Quantum Mechanics --- Lecture Notes}. University of Cologne.

\bibitem{boykin}
P.~O.~Boykin, M.~Sitharam, M.~Tarifi, and P.~Wocjan, ``Real Mutually Unbiased Bases,'' arXiv:quant-ph/0502024.

\bibitem{barandes}
J.~A.~Barandes, ``The Stochastic-Quantum Correspondence,'' arXiv:2302.10778.

\bibitem{araki1980}
H.~Araki, ``On a Characterization of the State Space of Quantum Mechanics,'' Communications in Mathematical Physics 75, 1--24 (1980).

\bibitem{hardy2001}
L.~Hardy, ``Quantum Theory From Five Reasonable Axioms,'' arXiv:quant-ph/0101012.

\bibitem{adler1995}
S.~L.~Adler, \textit{Quaternionic Quantum Mechanics and Quantum Fields}. Oxford University Press, 1995.

\bibitem{renou2021}
M.-O.~Renou, D.~Trillo, M.~Weilenmann, T.~P.~Le, A.~Tavakoli, N.~Gisin, A.~Ac\'in, and M.~Navascu\'es, ``Quantum theory based on real numbers can be experimentally falsified,'' Nature 600, 625--629 (2021).

\bibitem{hoffreumon-woods-2025}
T.~Hoffreumon and M.~P.~Woods, ``Quantum theory does not need complex numbers,'' arXiv:2504.02808 (2025).

\bibitem{barrios-hita-2026}
P.~Barrios Hita, A.~Trushechkin, H.~Kampermann, M.~Epping, and D.~Bru\ss, ``Quantum Mechanics Based on Real Numbers: A Consistent Description,'' Phys.\ Rev.\ Lett.\ 136, 240202 (2026).

\bibitem{needham1997}
T.~Needham, \textit{Visual Complex Analysis}. Clarendon Press / Oxford University Press, 1997.

\bibitem{zyczkowski2003}
K.~Życzkowski, M.~Kuś, W.~Słomczyński, and H.-J.~Sommers, ``Random unistochastic matrices,'' J.\ Phys.\ A: Math.\ Gen.\ 36, 3425--3450 (2003).

\end{thebibliography}

\printindex

\end{document}
